\chapter{Experimental Platform and Multi-Server Methodology}
\label{ch:methodology_b}

\section{Overview}

Thesis~A established that a single Dockerised microservice can exhibit realistic capacity-planning phenomena — growing mean latency, sharp tail-latency amplification, and eventual CPU saturation — on commodity hardware. Thesis~B expands that foundation into a systematic, multi-server experimental programme. Nine distinct server implementations across five languages and two concurrency configurations are subjected to open-loop load testing at a range of arrival rates, and the resulting traces are used both to challenge and to extend classical M/G/$c$ queueing assumptions.

The core research question that emerged from this expanded programme is:

\begin{quote}
\emph{Does the M/G/$c$ model's constant-service-time assumption hold for real containerised microservices, and if not, how can the model be corrected?}
\end{quote}

The answer — that mean service time grows with utilisation for CPU-bound servers — constitutes the central modelling contribution of this thesis.

%-------------------------------------------------------------
\section{Server Catalogue}
\label{sec:server-catalogue}

Nine server types were deployed, covering a representative cross-section of language runtimes, concurrency models, and workload characteristics. Table~\ref{tab:server-catalogue} summarises the configuration of each.

\begin{table}[h]
\centering
\caption{Server implementations used in Thesis~B experiments. Column $c$ is the number of concurrent workers modelled in the M/G/$c$ DES.}
\label{tab:server-catalogue}
\begin{tabular}{l l c l}
\toprule
\textbf{Tag} & \textbf{Language / Stack} & \textbf{$c$} & \textbf{Notes} \\
\midrule
Apache DSP (1c)   & PHP via Apache \texttt{mpm\_prefork} & 1 & One OS process per connection \\
Apache MSG (1c)   & PHP via Apache \texttt{mpm\_prefork} & 1 & Message-passing variant \\
Go lognormal (1c) & Go (goroutines)                     & 1 & CPU-bound arithmetic task \\
Go SQLite (1c)    & Go + SQLite                         & 1 & I/O-bound database reads \\
Java DSP (1c)     & Java (JVM, JIT)                     & 4 & Thread pool; JIT warmup present \\
Node DSP (1c)     & Node.js (single event loop)         & 1 & Non-blocking async handlers \\
Python DSP (1c)   & Python (GIL-constrained)            & 1 & Global interpreter lock enforced \\
Python DSP (3c)   & Python (3 worker processes)         & 3 & 3 independent processes \\
Node DSP (3c)     & Node.js cluster (3 workers)         & 3 & \texttt{cluster} module, Bernoulli splitting \\
\bottomrule
\end{tabular}
\end{table}

Each server was deployed as a Docker container with CPU resources pinned using both \texttt{--cpuset-cpus} and \texttt{--cpus} flags to enforce hard resource limits and prevent OS scheduling interference from other cores. This is critical: without \texttt{--cpuset-cpus}, the container scheduler may migrate threads across physical cores, masking throttling effects.

%-------------------------------------------------------------
\section{Load Generation and Trace Collection}
\label{sec:load-gen}

\subsection{Open-Loop Load Testing with Vegeta}

All experiments use Vegeta in open-loop mode: requests are issued at a fixed rate $\lambda$ (req/s) independently of server response times. This is the correct choice for studying queueing behaviour, because open-loop generation allows the queue to grow naturally under load, whereas closed-loop generation would cap the effective arrival rate and mask saturation.

For each server, experiments were run at a range of arrival rates from low load ($\rho \approx 0.1$) up to near-saturation ($\rho \approx 0.9$). Saturated traces ($\rho > 0.92$) were excluded from model fitting to avoid artefacts from dropped or rejected requests.

\subsection{Trace Format}

Each experiment produces a per-request CSV trace with the following columns:

\begin{itemize}
    \item \texttt{arrival\_unix\_ns} — absolute arrival timestamp in nanoseconds;
    \item \texttt{response\_ms} — end-to-end response time observed by the client;
    \item \texttt{service\_ms} — server-side processing time (measured within the handler using high-resolution timers where available);
    \item \texttt{queue\_ms} — inferred queueing delay (\texttt{response\_ms} $-$ \texttt{service\_ms});
    \item \texttt{status\_code} — HTTP status; non-2xx rows are excluded from analysis.
\end{itemize}

\subsection{Summary Statistics}

For each server $\times$ rate combination, a summary CSV records aggregate statistics computed from the trace: \texttt{rate\_rps}, \texttt{mean\_resp\_ms}, \texttt{mean\_svc\_ms}, \texttt{rho} ($= \lambda \cdot E[S] / (c \cdot 1000)$), and tail-latency percentiles. These summaries serve as the primary input to the service-time model-fitting step.

%-------------------------------------------------------------
\section{Concurrency Models}
\label{sec:concurrency}

A key design decision in applying M/G/$c$ is determining the correct value of $c$ — the number of parallel servers — for each implementation.

\paragraph{Apache \texttt{mpm\_prefork}.}
Each incoming HTTP connection spawns a dedicated OS process. When the container is pinned to a single CPU core, these processes compete for the core through OS preemption: the kernel time-slices among all active handler processes. This is a Processor Sharing (PS) discipline at the OS level, even though the application sees FCFS at the queue. The effective concurrency is $c = 1$ (one physical core), but the OS interleaves service among multiple simultaneous requests.

\paragraph{Go goroutines.}
The Go runtime multiplexes goroutines onto $\texttt{GOMAXPROCS}$ OS threads. With one pinned core, Go schedules goroutines cooperatively with preemption points, creating similar OS-level time-sharing to Apache but with lower context-switch overhead. $c = 1$.

\paragraph{Python GIL.}
The Global Interpreter Lock serialises Python bytecode execution: only one thread runs Python at a time. This effectively serialises service even when multiple threads are present. $c = 1$.

\paragraph{Node.js single process.}
The Node.js event loop is inherently single-threaded. $c = 1$. The 3-worker variant uses the \texttt{cluster} module, which forks three independent Node.js processes. The OS load-balances incoming connections across workers by Bernoulli splitting (roughly uniform), making this architecturally equivalent to three independent M/G/1 queues in parallel, \emph{not} a shared-queue M/G/3 system.

\paragraph{Java thread pool.}
The JVM maintains a thread pool of 4 worker threads. All threads share the CPU core, again creating OS-level time-sharing. $c = 4$ in the M/G/$c$ model.

The distinction between OS-level Processor Sharing (the physical reality) and application-level FCFS (the queueing abstraction) is central to understanding why service times grow with utilisation, as detailed in Chapter~\ref{ch:loaddep}.

%-------------------------------------------------------------
\section{Simulation Infrastructure}
\label{sec:des-infra}

\subsection{Batch Replay DES}

The core DES engine (\texttt{run\_des} in \texttt{analysis/des/des\_load\_dependent.py}) operates in \emph{replay mode}: it consumes the exact arrival sequence from an observed trace and produces a simulated response-time distribution by sampling service times from the fitted model.

\begin{enumerate}
    \item For each arrival at time $t_i$, assign it to the worker $s^* = \arg\min_s \texttt{free\_at}[s]$.
    \item Compute start time: $\texttt{start} = \max(t_i,\, \texttt{free\_at}[s^*])$.
    \item Sample service time: $\texttt{svc} \sim \text{Lognormal}(S_\text{eff}, \text{CV})$.
    \item Update $\texttt{free\_at}[s^*] = \texttt{start} + \texttt{svc}$.
    \item Record response time $\texttt{free\_at}[s^*] - t_i$.
\end{enumerate}

This replay approach removes arrival-process modelling error: both observed and simulated sequences share the same inter-arrival times, so any difference in the CDF is attributable solely to the service-time model. Five independent runs (differing only in the random seed for service-time sampling) are merged to reduce sampling variance.

\subsection{M0/M1/M2 Service-Time Model Comparison}

After the initial load-dependent DES results, the supervisor requested a more explicit comparison between the measured response-time CDF and the response-time CDFs produced by each service-time hypothesis. This became a separate validation pipeline implemented in \texttt{analysis/des/des\_m0\_m1\_m2\_compare.py}.

For each stable trace, the same arrival sequence is replayed through three DES variants:

\begin{itemize}
    \item \textbf{M0}: constant mean service time $E[S]=S_0$;
    \item \textbf{M1}: linear mean service time $E[S]=S_0+\alpha\rho_0$;
    \item \textbf{M2}: hyperbolic mean service time $E[S]=S_0/(1-\beta\rho_0)$.
\end{itemize}

The output table \texttt{results/tables/des\_m0\_m1\_m2\_results.csv} records the KS distance from the measured response-time CDF for each model. A follow-up script, \texttt{analysis/des/plot\_supervisor\_followup.py}, extracts the clearest cases where M2 improves on M0 and plots them in \texttt{results/figures/m2\_beats\_m0\_cdf.png}. This comparison is important because a service-time model should not only fit measured service times; it should also improve the simulated response-time distribution.

\subsection{Lognormal Service Time Parameterisation}

Service times are sampled from a Lognormal distribution parameterised by mean $\mu_S$ and coefficient of variation $\text{CV}$:

\[
\sigma^2 = \ln(1 + \text{CV}^2), \qquad
\mu = \ln\mu_S - \tfrac{\sigma^2}{2}.
\]

The CV is estimated once per server from the lowest-load trace (where queueing delay is negligible and the observed response-time distribution approximates the service-time distribution). It is held fixed across all load levels, capturing the shape of the service-time distribution while the mean $\mu_S$ varies with load.

\subsection{Validation Metric: KS Distance}

The Kolmogorov–Smirnov (KS) distance between two empirical distributions $F_1$ and $F_2$ is:

\[
d_\text{KS} = \sup_x |F_1(x) - F_2(x)|.
\]

A two-pointer $O(n + m)$ implementation operating on pre-sorted arrays is used. Lower KS distance indicates better agreement between the simulated and observed response-time CDFs. For the load-dependent work, the key comparisons are $d_\text{KS}(\text{CDF\_REF}, \text{CDF\_M0})$, $d_\text{KS}(\text{CDF\_REF}, \text{CDF\_M1})$, and $d_\text{KS}(\text{CDF\_REF}, \text{CDF\_M2})$.

\subsection{Service-Time Distribution Inspection and ANOVA}

The final analysis step tests the constant-service-time assumption directly at the trace level. The script \texttt{analysis/des/anova\_service\_time.py} groups per-request \texttt{service\_ms} samples by arrival rate and applies one-way ANOVA to $\log(\texttt{service\_ms})$. The log transform is used because the empirical service-time histograms are right-skewed and approximately log-normal.

The null hypothesis is that the mean log service time is identical across arrival rates, which corresponds to the M0 assumption that service time is independent of load. The alternative hypothesis is that service time depends on arrival rate. The output is \texttt{results/tables/service\_time\_anova.csv}; interpretation uses both $p$-value and $\eta^2$ effect size, because the large trace sizes make small differences statistically significant.

\subsection{Interactive Visualiser}

A browser-based visualiser (\texttt{visualiser.html}) allows interactive exploration of simulation results. Features include:

\begin{itemize}
    \item real-time DES with adjustable arrival rate $\lambda$, worker count $c$, mean service time, and CV;
    \item CSV trace loader: uploading an experiment trace overlays the observed CDF and displays the KS distance;
    \item Gantt timeline panel: replays the loaded trace with server assignment shown as coloured bars, with a DES shadow for comparison;
    \item live experiment mode: a WebSocket backend (\texttt{vis\_server.py}) fires real HTTP requests at any running Docker server and streams per-request timings into the browser in real time.
\end{itemize}
